% =====================================================================
\chapter{Derivadas}

No capítulo anterior usamos limites para calcular a inclinação da reta tangente e a
velocidade instantânea. Essas duas ideias são casos de um único conceito: a
\emph{derivada}, que mede a taxa de variação instantânea de uma função. Neste capítulo
definimos a derivada, aprendemos regras para calculá-la sem recorrer ao limite a cada vez e
a aplicamos a movimento, crescimento populacional, custo e receita. Ao final teremos as
derivadas de todas as funções do capítulo~1.

% ---------------------------------------------------------------------
\section{Definindo a derivada}

\subsection{Retas tangentes}

Seja $f$ definida em um intervalo $I$ que contém $a$. Para $x\neq a$ em $I$, o
\textbf{quociente de diferenças}\index{quociente de diferencas@quociente de diferenças}
\[
Q=\frac{f(x)-f(a)}{x-a}
\qquad\text{ou, com } x=a+h,\qquad
Q=\frac{f(a+h)-f(a)}{h}
\]
é a inclinação da reta secante pelos pontos $(a,f(a))$ e $(x,f(x))$. Quando $x\to a$
(ou $h\to0$), as secantes se aproximam da reta tangente.

\figtikz{\begin{tikzpicture}
\begin{axis}[painel,width=8cm,height=6cm,xmin=-0.3,xmax=3.2,ymin=-1,ymax=6.5,xtick={1},xticklabels={$a$},ytick=\empty]
\addplot[domain=-0.3:2.5]{x^2};
\foreach \h/\o in {1.4/30,0.8/50,0.4/70}{
  \edef\tmp{\noexpand\addplot[laranja!\o!white,thin,domain=-0.3:3.2]{1+(2+\h)*(x-1)};
  \noexpand\fill[laranja!\o!white] (axis cs:{1+\h},{(1+\h)^2}) circle (1.6pt);}\tmp}
\addplot[laranja,thick,domain=-0.3:3.2]{1+2*(x-1)};
\fill[black] (axis cs:1,1) circle (2pt) node[below right,font=\footnotesize]{$(a,f(a))$};
\node[laranja,font=\footnotesize] at (axis cs:2.75,4.1) {tangente};
\end{axis}\end{tikzpicture}}{As secantes por $(a,f(a))$ e por pontos cada vez mais
próximos (laranja-claro) tendem à reta tangente (laranja).}

\begin{definicao}
Seja $f$ definida em um intervalo aberto contendo $a$. A \textbf{reta tangente} ao gráfico
de $f$ em $a$ é a reta que passa por $(a,f(a))$ com inclinação
\[
m_{\tan}=\lim_{x\to a}\frac{f(x)-f(a)}{x-a}=\lim_{h\to0}\frac{f(a+h)-f(a)}{h},
\]
desde que esse limite exista.
\end{definicao}

\begin{exemplo}{3.1}{Reta tangente}
Encontre a reta tangente a $f(x)=x^2$ em $x=3$ e a $g(x)=1/x$ em $x=2$.
\solucao
$m_{\tan}=\lim_{x\to3}\dfrac{x^2-9}{x-3}=\lim_{x\to3}(x+3)=6$; a tangente é
$y-9=6(x-3)$, ou $y=6x-9$.

Para $g$: $m_{\tan}=\lim_{h\to0}\dfrac{\frac1{2+h}-\frac12}{h}
=\lim_{h\to0}\dfrac{-h}{2(2+h)h}=-\dfrac14$; a tangente é $y-\tfrac12=-\tfrac14(x-2)$, ou
$y=-\tfrac14x+1$.
\end{exemplo}

\subsection{A derivada em um ponto}

A inclinação da tangente é tão importante que recebe um nome próprio.

\begin{definicao}
A \textbf{derivada}\index{derivada} de $f$ em $a$, denotada $f'(a)$, é
\[
f'(a)=\lim_{x\to a}\frac{f(x)-f(a)}{x-a}=\lim_{h\to0}\frac{f(a+h)-f(a)}{h},
\]
desde que o limite exista. Nesse caso, dizemos que $f$ é \textbf{derivável}\index{funcao@função!derivavel@derivável} (ou
diferenciável) em $a$. A derivada $f'(a)$ é a \textbf{taxa de variação instantânea}\index{taxa de variacao@taxa de variação} de $f$
em $a$.
\end{definicao}

\begin{exemplo}{3.2}{Calculando uma derivada pela definição}
Para $f(x)=3x^2-4x+1$, calcule $f'(2)$.
\solucao
\begin{align*}
f'(2)&=\lim_{h\to0}\frac{\big(3(2+h)^2-4(2+h)+1\big)-5}{h}
=\lim_{h\to0}\frac{12+12h+3h^2-8-4h+1-5}{h}\\
&=\lim_{h\to0}\frac{8h+3h^2}{h}=\lim_{h\to0}(8+3h)=8.
\end{align*}
\end{exemplo}

\subsection{Velocidade e taxas de variação}

Se $s(t)$ é a posição de um objeto que se move ao longo de uma reta, a velocidade média em
$[a,a+h]$ é $\frac{s(a+h)-s(a)}{h}$, e a \textbf{velocidade instantânea} em $t=a$ é o
limite dessas médias: $v(a)=s'(a)$. Mais geralmente, para qualquer grandeza $f$, a taxa
média de variação em um intervalo é um quociente de diferenças, e a taxa instantânea é a
derivada.

\begin{exemplo}{3.3}{Velocidade de uma bola em queda}
Uma bola solta do alto de um prédio de 64 pés tem altura $s(t)=-16t^2+64$ (em pés, com $t$
em segundos). Qual é sua velocidade em $t=1$?
\solucao
$s'(1)=\lim_{t\to1}\dfrac{-16t^2+64-48}{t-1}=\lim_{t\to1}\dfrac{-16(t^2-1)}{t-1}
=\lim_{t\to1}-16(t+1)=-32$ pés/s. O sinal negativo indica que a bola está descendo.
\end{exemplo}

% ---------------------------------------------------------------------
\section{A derivada como função}

\begin{definicao}
A \textbf{função derivada} $f'$ é a função cujo domínio é formado pelos $x$ para os quais
existe
\[
f'(x)=\lim_{h\to0}\frac{f(x+h)-f(x)}{h}.
\]
Também escrevemos $\dfrac{dy}{dx}$, $\dfrac{df}{dx}$, $\dfrac{d}{dx}\big(f(x)\big)$ ou $y'$
(\emph{notação de Leibniz} e \emph{de Lagrange}).
\end{definicao}

\begin{exemplo}{3.4}{Derivada da raiz quadrada}
Para $f(x)=\sqrt x$:
\[
f'(x)=\lim_{h\to0}\frac{\sqrt{x+h}-\sqrt x}{h}\cdot\frac{\sqrt{x+h}+\sqrt x}{\sqrt{x+h}+\sqrt x}
=\lim_{h\to0}\frac{h}{h(\sqrt{x+h}+\sqrt x)}=\frac{1}{2\sqrt x}.
\]
O domínio de $f'$ é $(0,\infty)$, menor que o de $f$: em $x=0$ a tangente é vertical.
\end{exemplo}

\subsection{O gráfico da derivada}

O valor de $f'(x)$ é a inclinação do gráfico de $f$ em $x$. Assim:
\begin{itemize}
\item onde $f$ é crescente, $f'>0$; onde $f$ é decrescente, $f'<0$;
\item onde o gráfico de $f$ tem tangente horizontal, $f'=0$.
\end{itemize}

\figtikz{\begin{tikzpicture}
\begin{groupplot}[group style={group size=1 by 2,vertical sep=1.1cm},painel,width=9cm,height=4.4cm,
  xmin=-2.4,xmax=2.4,xtick={-1,1},ytick={-2,2}]
\nextgroupplot[ymin=-3,ymax=3,title={$f(x)=x^3-3x$}]
\addplot[domain=-2.4:2.4]{x^3-3*x};
\draw[laranja,thick] (axis cs:-1.6,2) -- (axis cs:-0.4,2);
\draw[laranja,thick] (axis cs:0.4,-2) -- (axis cs:1.6,-2);
\draw[cinza,dotted] (axis cs:-1,-3) -- (axis cs:-1,3); \draw[cinza,dotted] (axis cs:1,-3) -- (axis cs:1,3);
\nextgroupplot[ymin=-4,ymax=8,title={$f'(x)=3x^2-3$}]
\addplot[domain=-2.4:2.4,laranja]{3*x^2-3};
\draw[cinza,dotted] (axis cs:-1,-4) -- (axis cs:-1,8); \draw[cinza,dotted] (axis cs:1,-4) -- (axis cs:1,8);
\end{groupplot}\end{tikzpicture}}{Onde $f$ cresce ($|x|>1$), $f'>0$; onde $f$ decresce
($|x|<1$), $f'<0$; nos pontos de tangente horizontal ($x=\pm1$), $f'=0$.}

\subsection{Derivabilidade e continuidade}

\begin{teorema}[Derivabilidade implica continuidade]{3.1}\index{derivabilidade e continuidade}
Se $f$ é derivável em $a$, então $f$ é contínua em $a$.
\end{teorema}
\emph{Prova.} Para $x\neq a$,
$f(x)-f(a)=\dfrac{f(x)-f(a)}{x-a}\,(x-a)\ \to\ f'(a)\cdot0=0$ quando $x\to a$.
Logo $\lim_{x\to a}f(x)=f(a)$. $\square$

A recíproca é falsa: uma função pode ser contínua em $a$ e não ser derivável ali. Há três
situações típicas:

\figtikz{\begin{tikzpicture}
\begin{groupplot}[group style={group size=3 by 1,horizontal sep=1.0cm},painel,width=5.3cm,height=4.4cm,
  xmin=-1.5,xmax=1.5,xtick=\empty,ytick=\empty]
\nextgroupplot[title={(a) bico: $|x|$},ymin=-0.4,ymax=1.5]
\addplot[domain=-1.5:1.5]{abs(x)};
\nextgroupplot[title={(b) tangente vertical: $\sqrt[3]{x}$},ymin=-1.3,ymax=1.3]
\addplot[domain=0:1.5,samples=200]{x^(1/3)}; \addplot[domain=-1.5:0,samples=200]{-(-x)^(1/3)};
\draw[laranja,dashed] (axis cs:0,-1.3) -- (axis cs:0,1.3);
\nextgroupplot[title={(c) oscilação: $x\sin\frac1x$},ymin=-0.5,ymax=0.5,xmin=-0.5,xmax=0.5]
\addplot[domain=-0.5:-0.002,samples=1500,thin]{x*sin(deg(1/x))};
\addplot[domain=0.002:0.5,samples=1500,thin]{x*sin(deg(1/x))};
\addplot[laranja,dashed,thin,domain=-0.5:0.5]{x}; \addplot[laranja,dashed,thin,domain=-0.5:0.5]{-x};
\end{groupplot}\end{tikzpicture}}{Funções contínuas em $0$, mas não deriváveis em $0$.
(a) Os limites laterais do quociente de diferenças são $-1$ e $1$. (b) O quociente tende a
$+\infty$. (c) O quociente $\sin(1/h)$ oscila e não tem limite.}

\begin{exemplo}{3.5}{Uma função por partes derivável}
Encontre $a$ e $b$ para que $f(x)=\begin{cases}ax+b,&x<3\\x^2,&x\ge3\end{cases}$ seja
derivável em todo $\R$.
\solucao
Contínua em 3: $3a+b=9$. Para a derivada existir em 3, os quocientes laterais devem ter o
mesmo limite. À direita, a inclinação de $x^2$ em 3 é 6; à esquerda, é $a$. Logo $a=6$ e
$b=-9$. A reta $y=6x-9$ é exatamente a tangente à parábola em $x=3$.
\end{exemplo}

\subsection{Derivadas de ordem superior}

A derivada de $f'$ é a \textbf{segunda derivada}\index{derivada!de ordem superior}, $f''$ ou $\dfrac{d^2y}{dx^2}$; depois vêm
$f'''$, $f^{(4)},\dots$ Se $s(t)$ é a posição, $v(t)=s'(t)$ é a velocidade e
$a(t)=v'(t)=s''(t)$ é a \textbf{aceleração}\index{aceleracao@aceleração}.

% ---------------------------------------------------------------------
\section{Regras de derivação}

Calcular derivadas pela definição é trabalhoso. As regras a seguir, provadas a partir da
definição, tornam o cálculo mecânico.

\begin{teorema}[Constante e potência]{3.2}
Se $c$ é constante, $\dfrac{d}{dx}(c)=0$. Para todo inteiro $n\ge1$,
$\dfrac{d}{dx}\,x^n=nx^{n-1}$.
\end{teorema}
\emph{Prova da regra da potência.} Pelo binômio de Newton,
$(x+h)^n=x^n+nx^{n-1}h+\binom n2x^{n-2}h^2+\dots+h^n$. Então
\[
\frac{(x+h)^n-x^n}{h}=nx^{n-1}+\binom n2x^{n-2}h+\dots+h^{n-1}\ \longrightarrow\ nx^{n-1}. \qquad\square
\]

\begin{teorema}[Soma, diferença e múltiplo constante]{3.3}
Se $f$ e $g$ são deriváveis e $k$ é constante:
\[
(f\pm g)'=f'\pm g',\qquad (kf)'=k\,f'.
\]
\end{teorema}

Com isso, derivamos qualquer polinômio termo a termo: por exemplo,
$\dfrac{d}{dx}\big(2x^5+7x^2-x+4\big)=10x^4+14x-1$.

\begin{teorema}[Regra do produto]{3.4}\index{regra do produto}
Se $f$ e $g$ são deriváveis, $\big(f\,g\big)'=f'\,g+f\,g'$.
\end{teorema}
\emph{Prova.} Somamos e subtraímos $f(x+h)g(x)$ no numerador:
\[
\frac{f(x+h)g(x+h)-f(x)g(x)}{h}
=f(x+h)\,\frac{g(x+h)-g(x)}{h}+g(x)\,\frac{f(x+h)-f(x)}{h}
\ \longrightarrow\ f(x)g'(x)+g(x)f'(x),
\]
usando que $f(x+h)\to f(x)$, pois $f$ é derivável, logo contínua. $\square$

\begin{nota}
A derivada do produto \emph{não} é o produto das derivadas: $(x\cdot x)'=2x$, e não
$1\cdot1$.
\end{nota}

\begin{teorema}[Regra do quociente]{3.5}\index{regra do quociente}
Se $f$ e $g$ são deriváveis e $g(x)\neq0$,
\[
\Big(\frac fg\Big)'=\frac{f'\,g-f\,g'}{g^2}.
\]
\end{teorema}

Aplicando a regra do quociente a $x^{-n}=1/x^n$, obtemos
$\dfrac{d}{dx}x^{-n}=\dfrac{0-nx^{n-1}}{x^{2n}}=-nx^{-n-1}$: a regra da potência vale também
para expoentes inteiros negativos.

\begin{exemplo}{3.6}{Produto e quociente}
(a) $\dfrac{d}{dx}\big[(x^2+2)(3x^3-5x)\big]=2x(3x^3-5x)+(x^2+2)(9x^2-5)=15x^4+3x^2-10$.

(b) $\dfrac{d}{dx}\Big[\dfrac{5x^2}{4x+3}\Big]=\dfrac{10x(4x+3)-5x^2\cdot4}{(4x+3)^2}
=\dfrac{20x^2+30x}{(4x+3)^2}$.

(c) $\dfrac{d}{dx}\Big[\dfrac{6}{x^2}\Big]=\dfrac{d}{dx}\big[6x^{-2}\big]=-12x^{-3}=-\dfrac{12}{x^3}$.
\end{exemplo}

\begin{exemplo}{3.7}{Tangentes horizontais}
Onde o gráfico de $f(x)=x^3-6x^2+9x+1$ tem tangente horizontal?
\solucao
$f'(x)=3x^2-12x+9=3(x-1)(x-3)=0$ em $x=1$ e $x=3$. Os pontos são $(1,5)$ e $(3,1)$.
\end{exemplo}

% ---------------------------------------------------------------------
\section{Derivadas como taxas de variação}

\textbf{Fórmula da variação aproximada.} Como $f'(a)\approx\frac{f(a+h)-f(a)}h$ para $h$
pequeno,
\[
f(a+h)\approx f(a)+f'(a)\,h .
\]
Por exemplo, se $f(3)=2$ e $f'(3)=5$, então $f(3{,}2)\approx2+5(0{,}2)=3$.

\subsection{Movimento retilíneo}

\begin{definicao}
Se $s(t)$ é a posição de um objeto no instante $t$: a \textbf{velocidade} é
$v(t)=s'(t)$; a \textbf{rapidez} (módulo da velocidade) é $|v(t)|$; a
\textbf{aceleração} é $a(t)=v'(t)=s''(t)$.
\end{definicao}

O objeto se move no sentido positivo quando $v>0$, no negativo quando $v<0$, e está
momentaneamente em repouso quando $v=0$. Ele acelera (a rapidez aumenta) quando $v$ e $a$
têm o mesmo sinal, e desacelera quando têm sinais opostos.

\begin{exemplo}{3.8}{Posição e velocidade}
Uma partícula tem posição $s(t)=t^3-9t^2+24t$ ($t\ge0$). Quando ela está em repouso? Em
que intervalos se move para a frente e para trás?
\solucao
$v(t)=3t^2-18t+24=3(t-2)(t-4)$. Repouso em $t=2$ e $t=4$. Pelo sinal de $v$: para a frente
em $[0,2)$ e $(4,\infty)$, para trás em $(2,4)$. A aceleração é $a(t)=6t-18$, nula em $t=3$.
\end{exemplo}

\subsection{Populações, custo e receita}

Se $P(t)$ é uma população, $P'(t)$ é a \textbf{taxa de crescimento populacional}. Na
economia, se $C(x)$ é o custo de produzir $x$ unidades e $R(x)$ é a receita com a venda de
$x$ unidades, o \textbf{custo marginal}\index{analise marginal@análise marginal} é $C'(x)$, a \textbf{receita marginal} é $R'(x)$
e o \textbf{lucro marginal} é $P'(x)=R'(x)-C'(x)$. Pela fórmula da variação aproximada, o
custo marginal $C'(x)$ estima o custo de produzir a unidade número $x+1$.

\begin{exemplo}{3.9}{Receita marginal}
Se $R(x)=4x-0{,}01x^2$ (em mil reais), então $R'(x)=4-0{,}02x$. Para $x=100$, $R'(100)=2$:
vender a 101ª unidade aumenta a receita em cerca de R\$\,2 mil. O valor exato é
$R(101)-R(100)=1{,}99$.
\end{exemplo}

% ---------------------------------------------------------------------
\section{Derivadas das funções trigonométricas}

\begin{teorema}[Derivadas de seno e cosseno]{3.6}
\[
\frac{d}{dx}(\sin x)=\cos x,\qquad \frac{d}{dx}(\cos x)=-\sin x .
\]
\end{teorema}
\emph{Prova para o seno.} Pela fórmula da soma,
\[
\frac{\sin(x+h)-\sin x}{h}=\sin x\cdot\frac{\cos h-1}{h}+\cos x\cdot\frac{\sin h}{h}
\ \longrightarrow\ \sin x\cdot0+\cos x\cdot1=\cos x,
\]
usando os dois limites trigonométricos do Teorema~2.7. A prova para o cosseno é
análoga. $\square$

\figtikz{\begin{tikzpicture}
\begin{axis}[painel,width=11cm,height=4.4cm,xmin=-6.6,xmax=6.6,ymin=-1.5,ymax=1.5,
  xtick={-6.2832,-3.1416,3.1416,6.2832},xticklabels={$-2\pi$,$-\pi$,$\pi$,$2\pi$},ytick={-1,1},
  legend style={font=\footnotesize,at={(1,1.05)},anchor=south east},legend columns=2]
\addplot[domain=-6.6:6.6]{sin(deg(x))}; \addlegendentry{$y=\sin x$}
\addplot[laranja,domain=-6.6:6.6]{cos(deg(x))}; \addlegendentry{$y=\cos x$ (derivada)}
\foreach \p in {-4.7124,-1.5708,1.5708,4.7124}{\edef\tmp{\noexpand\fill[black] (axis cs:\p,0) circle (1.6pt);}\tmp}
\end{axis}\end{tikzpicture}}{Onde $\sin x$ tem tangente horizontal (máximos e mínimos), sua
derivada $\cos x$ se anula (pontos pretos); onde $\sin x$ cresce mais rápido ($x=0$), $\cos x$
vale 1.}

Pela regra do quociente, $\dfrac{d}{dx}\tan x=\dfrac{\cos x\cos x-\sin x(-\sin x)}{\cos^2x}
=\dfrac1{\cos^2x}=\sec^2x$. Do mesmo modo:

\begin{teorema}[Demais funções trigonométricas]{3.7}
\[
(\tan x)'=\sec^2x,\qquad (\cot x)'=-\csc^2x,\qquad (\sec x)'=\sec x\tan x,\qquad
(\csc x)'=-\csc x\cot x .
\]
\end{teorema}

\textbf{Derivadas de ordem superior.} As derivadas de $\sin x$ se repetem com período 4:
$\sin x\to\cos x\to-\sin x\to-\cos x\to\sin x$. Assim,
$\dfrac{d^{74}}{dx^{74}}\sin x=\dfrac{d^2}{dx^2}\sin x=-\sin x$, pois $74=4\cdot18+2$.

\begin{exemplo}{3.10}{Uma massa presa a uma mola}
A posição de uma massa é $s(t)=2\sin t$ (cm). Então $v(t)=2\cos t$ e $a(t)=-2\sin t=-s(t)$:
a aceleração é proporcional ao deslocamento e tem sentido oposto. Em $t=\pi/2$, a massa está
em $s=2$, parada ($v=0$), com aceleração $-2$ cm/s$^2$, que a puxa de volta.
\end{exemplo}

% ---------------------------------------------------------------------
\section{A regra da cadeia}

Como derivar $h(x)=\sin(x^3)$ ou $(3x^2+1)^{20}$? São funções compostas, e sua derivada
segue uma regra simples: se $u=g(x)$ varia $g'(x)$ vezes mais rápido que $x$, e $y=f(u)$
varia $f'(u)$ vezes mais rápido que $u$, então $y$ varia $f'(u)\,g'(x)$ vezes mais rápido que
$x$.

\begin{regra}{Regra da cadeia}\index{regra da cadeia}
Se $g$ é derivável em $x$ e $f$ é derivável em $g(x)$, então $h=f\circ g$ é derivável em
$x$ e
\[
h'(x)=f'\big(g(x)\big)\,g'(x).
\]
Na notação de Leibniz, com $y=f(u)$ e $u=g(x)$:\qquad
$\dfrac{dy}{dx}=\dfrac{dy}{du}\cdot\dfrac{du}{dx}$.
\end{regra}

Em palavras: \emph{derive a função de fora, mantendo a de dentro, e multiplique pela
derivada da de dentro}. Dois casos frequentes:
\[
\frac{d}{dx}\big(g(x)\big)^n=n\big(g(x)\big)^{n-1}g'(x),\qquad
\frac{d}{dx}\sin\big(g(x)\big)=\cos\big(g(x)\big)\,g'(x),
\]
e analogamente $\frac{d}{dx}\cos u=-\sin u\,u'$, $\frac{d}{dx}\tan u=\sec^2u\,u'$, etc.

\begin{exemplo}{3.11}{Usando a regra da cadeia}
(a) $\dfrac{d}{dx}(3x^2+1)^3=3(3x^2+1)^2\cdot6x=18x(3x^2+1)^2$.

(b) $\dfrac{d}{dx}\sin^3x=3\sin^2x\cos x$.

(c) $\dfrac{d}{dx}\cos(5x^2)=-\sin(5x^2)\cdot10x$.

(d) $\dfrac{d}{dx}\Big[\dfrac{1}{(3x^2+1)^2}\Big]=\dfrac{d}{dx}(3x^2+1)^{-2}=-2(3x^2+1)^{-3}\cdot6x
=-\dfrac{12x}{(3x^2+1)^3}$.
\end{exemplo}

\textbf{Composição de três funções.} Para $k(x)=h\big(f(g(x))\big)$, aplicamos a regra duas
vezes: $k'(x)=h'\big(f(g(x))\big)\cdot f'\big(g(x)\big)\cdot g'(x)$. Por exemplo,
\[
\frac{d}{dx}\sin^4(7x-3)=4\sin^3(7x-3)\cdot\cos(7x-3)\cdot7 .
\]

\begin{exemplo}{3.12}{Cadeia e produto juntos}
$\dfrac{d}{dx}\big[x^2\cos(3x)\big]=2x\cos(3x)+x^2\cdot\big(-\sin(3x)\cdot3\big)
=2x\cos(3x)-3x^2\sin(3x)$.
\end{exemplo}

% ---------------------------------------------------------------------
\section{Derivadas de funções inversas}

Se $(a,b)$ está no gráfico de $f$, então $(b,a)$ está no gráfico de $f^{-1}$, que é a
reflexão em torno de $y=x$. A reflexão troca $\Delta x$ e $\Delta y$; logo, as inclinações
das tangentes em pontos correspondentes são \emph{recíprocas}.

\figtikz{\begin{tikzpicture}
\begin{axis}[painel,width=7cm,height=7cm,xmin=-2.2,xmax=4.2,ymin=-2.2,ymax=4.2,axis equal image,xtick=\empty,ytick=\empty]
\addplot[cinza,dotted,domain=-2.2:4.2]{x};
\addplot[domain=-2.2:1.45]{exp(x)}; \addplot[laranja,domain=0.12:4.2]{ln(x)};
\addplot[azul,dashed,thin,domain=-0.6:2.2]{exp(0.8)+exp(0.8)*(x-0.8)};
\addplot[laranja,dashed,thin,domain=0.6:4.2]{0.8+(x-exp(0.8))/exp(0.8)};
\fill[azul] (axis cs:0.8,{exp(0.8)}) circle (2pt) node[left,font=\footnotesize]{$(a,b)$};
\fill[laranja] (axis cs:{exp(0.8)},0.8) circle (2pt) node[below right,font=\footnotesize]{$(b,a)$};
\node[azul,font=\footnotesize] at (axis cs:-1.2,1.4) {$y=e^x$};
\node[laranja,font=\footnotesize] at (axis cs:3.4,0.7) {$y=\ln x$};
\end{axis}\end{tikzpicture}}{As tangentes a $f$ em $(a,b)$ e a $f^{-1}$ em $(b,a)$ são
reflexões uma da outra; suas inclinações são $m$ e $1/m$.}

\begin{teorema}[Derivada da função inversa]{3.8}
Seja $f$ inversível e derivável. Se $f'\big(f^{-1}(x)\big)\neq0$, então $f^{-1}$ é derivável em
$x$ e
\[
\big(f^{-1}\big)'(x)=\frac{1}{f'\big(f^{-1}(x)\big)} .
\]
\end{teorema}

\begin{exemplo}{3.13}{Aplicando o teorema da inversa}
Para $f(x)=x^3+x+1$, que é crescente, calcule $(f^{-1})'(3)$.
\solucao
$f(1)=3$, logo $f^{-1}(3)=1$. Como $f'(x)=3x^2+1$, temos $f'(1)=4$ e
$(f^{-1})'(3)=\dfrac1{f'(1)}=\dfrac14$. Não foi preciso encontrar a fórmula de $f^{-1}$.
\end{exemplo}

\textbf{Expoentes racionais.} Aplicando o teorema a $g(x)=x^{1/n}$, inversa de $x^n$,
obtemos $g'(x)=\frac1n x^{1/n-1}$. Combinando com a regra da cadeia, para todo racional $r$:
\[
\frac{d}{dx}x^r=r\,x^{r-1}.
\]
Por exemplo, $\dfrac{d}{dx}\sqrt[3]{x^2}=\dfrac{d}{dx}x^{2/3}=\dfrac23x^{-1/3}$.

\textbf{Inversas trigonométricas.} Para $y=\sin^{-1}x$, temos $\sin y=x$ com
$-\frac\pi2\le y\le\frac\pi2$. Pelo teorema, $\dfrac{dy}{dx}=\dfrac1{\cos y}$. Como
$\cos y\ge0$ nesse intervalo, $\cos y=\sqrt{1-\sin^2y}=\sqrt{1-x^2}$. Do mesmo modo:

\begin{teorema}[Derivadas das inversas trigonométricas]{3.9}
\[
\frac{d}{dx}\sin^{-1}x=\frac{1}{\sqrt{1-x^2}},\qquad
\frac{d}{dx}\cos^{-1}x=\frac{-1}{\sqrt{1-x^2}},\qquad
\frac{d}{dx}\tan^{-1}x=\frac{1}{1+x^2},
\]
\[
\frac{d}{dx}\cot^{-1}x=\frac{-1}{1+x^2},\qquad
\frac{d}{dx}\sec^{-1}x=\frac{1}{|x|\sqrt{x^2-1}},\qquad
\frac{d}{dx}\csc^{-1}x=\frac{-1}{|x|\sqrt{x^2-1}} .
\]
\end{teorema}

\begin{exemplo}{3.14}{Cadeia com inversa trigonométrica}
$\dfrac{d}{dx}\tan^{-1}(x^2)=\dfrac{2x}{1+x^4}$\qquad e\qquad
$\dfrac{d}{dx}\big[x\sin^{-1}x\big]=\sin^{-1}x+\dfrac{x}{\sqrt{1-x^2}}$.
\end{exemplo}

% ---------------------------------------------------------------------
\section{Derivação implícita}

Muitas curvas, como a circunferência $x^2+y^2=25$, não são gráficos de uma única função
$y=f(x)$. Ainda assim, perto de cada ponto (salvo exceções) a curva é o gráfico de alguma
função $y(x)$, e podemos derivar a equação \emph{como está}, tratando $y$ como função de $x$.

\begin{estrategia}{derivação implícita}\index{derivacao implicita@derivação implícita}
\begin{enumerate}[label=\arabic*.]
\item Derive os dois lados da equação em relação a $x$, lembrando que, pela regra da
  cadeia, $\frac{d}{dx}\big(g(y)\big)=g'(y)\,\frac{dy}{dx}$ (por exemplo,
  $\frac{d}{dx}y^2=2y\,y'$).
\item Passe os termos com $\frac{dy}{dx}$ para um lado e os demais para o outro.
\item Ponha $\frac{dy}{dx}$ em evidência e isole.
\end{enumerate}
\end{estrategia}

\begin{exemplo}{3.15}{Tangente a uma circunferência}
Na circunferência $x^2+y^2=25$: $2x+2y\,y'=0$, logo $y'=-\dfrac xy$. No ponto $(3,-4)$,
$y'=\frac34$, e a tangente é $y+4=\tfrac34(x-3)$. Note que a tangente é perpendicular ao raio,
cuja inclinação é $-\tfrac43$.
\end{exemplo}

\begin{exemplo}{3.16}{O fólio de Descartes}
Encontre a tangente à curva $x^3+y^3=6xy$ no ponto $(3,3)$.
\solucao
Derivando: $3x^2+3y^2y'=6y+6x\,y'$. Isolando,
\[
y'=\frac{2y-x^2}{y^2-2x}.
\]
Em $(3,3)$: $y'=\frac{6-9}{9-6}=-1$. A tangente é $y=-x+6$.
\figtikz{\begin{tikzpicture}
\begin{axis}[painel,width=6.2cm,height=6.2cm,xmin=-4,xmax=5,ymin=-4,ymax=5,axis equal image,xtick={3},ytick={3}]
\addplot[domain=-0.55:12,samples=400,variable=t] ({6*t/(1+t^3)},{6*t^2/(1+t^3)});
\addplot[domain=-50:-1.6,samples=300,variable=t] ({6*t/(1+t^3)},{6*t^2/(1+t^3)});
\addplot[laranja,domain=0.5:5.5]{-x+6};
\fill[laranja] (axis cs:3,3) circle (2pt);
\end{axis}\end{tikzpicture}}{O fólio de Descartes e a tangente em $(3,3)$.}
\end{exemplo}

% ---------------------------------------------------------------------
\section{Derivadas das funções exponenciais e logarítmicas}

Vimos no capítulo~1 que $e$ é a base cuja exponencial tem tangente de inclinação 1 em
$x=0$: $\lim_{h\to0}\frac{e^h-1}{h}=1$. Daí
\[
\frac{e^{x+h}-e^x}{h}=e^x\cdot\frac{e^h-1}{h}\ \longrightarrow\ e^x .
\]

\begin{teorema}[Derivada da exponencial natural]{3.10}
$\dfrac{d}{dx}e^x=e^x$, e, pela regra da cadeia, $\dfrac{d}{dx}e^{g(x)}=e^{g(x)}\,g'(x)$.
\end{teorema}

Como $\ln x$ é a inversa de $e^x$, o teorema da inversa dá
$\frac{d}{dx}\ln x=\frac{1}{e^{\ln x}}=\frac1x$.

\begin{teorema}[Derivada do logaritmo natural]{3.11}
Para $x>0$, $\dfrac{d}{dx}\ln x=\dfrac1x$, e $\dfrac{d}{dx}\ln g(x)=\dfrac{g'(x)}{g(x)}$
quando $g(x)>0$. Mais geralmente, $\dfrac{d}{dx}\ln|x|=\dfrac1x$ para $x\neq0$.
\end{teorema}

Para outras bases, escrevemos $b^x=e^{x\ln b}$ e $\log_bx=\frac{\ln x}{\ln b}$:

\begin{teorema}[Bases quaisquer]{3.12}
Para $b>0$, $b\neq1$:\qquad
$\dfrac{d}{dx}b^x=b^x\ln b$,\qquad $\dfrac{d}{dx}\log_bx=\dfrac{1}{x\ln b}$.
\end{teorema}

\begin{exemplo}{3.17}{Exponenciais e logaritmos}
(a) $\dfrac{d}{dx}\Big[\dfrac{e^{x^2}}{x}\Big]=\dfrac{2xe^{x^2}\cdot x-e^{x^2}}{x^2}
=\dfrac{e^{x^2}(2x^2-1)}{x^2}$.

(b) $\dfrac{d}{dx}\ln(x^3+3x-4)=\dfrac{3x^2+3}{x^3+3x-4}$.

(c) $\dfrac{d}{dx}\ln\Big(\dfrac{x^2\sin x}{2x+1}\Big)$: primeiro usamos as propriedades do
logaritmo, $2\ln x+\ln\sin x-\ln(2x+1)$; a derivada é $\dfrac2x+\cot x-\dfrac{2}{2x+1}$.

(d) $\dfrac{d}{dx}\,3^{x^2}=3^{x^2}\ln3\cdot2x$.
\end{exemplo}

\subsection{Derivação logarítmica}

Para funções da forma $f(x)^{g(x)}$, ou produtos e quocientes complicados, tomamos o
logaritmo dos dois lados antes de derivar.

\begin{estrategia}{derivação logarítmica}\index{derivacao logaritmica@derivação logarítmica}
\begin{enumerate}[label=\arabic*.]
\item Aplique $\ln$ aos dois lados de $y=h(x)$ e simplifique com as propriedades do
  logaritmo.
\item Derive implicitamente: o lado esquerdo vira $\frac{y'}{y}$.
\item Multiplique por $y$ e substitua $y=h(x)$.
\end{enumerate}
\end{estrategia}

\begin{exemplo}{3.18}{A derivada de $x^x$}
Seja $y=x^x$ ($x>0$). Então $\ln y=x\ln x$, e derivando,
$\dfrac{y'}{y}=\ln x+x\cdot\dfrac1x=\ln x+1$. Logo
\[
\frac{d}{dx}x^x=x^x(\ln x+1).
\]
Atenção: $x^x$ não é potência ($x^n$) nem exponencial ($b^x$), e nenhuma das duas regras
se aplica diretamente.
\end{exemplo}

Com a mesma técnica, para qualquer real $r$ e $x>0$, $y=x^r$ dá $\ln y=r\ln x$ e
$y'=y\cdot\frac rx=rx^{r-1}$: a regra da potência vale para \emph{todo} expoente real.

% ---------------------------------------------------------------------
\section*{Resumo do capítulo}\markright{Resumo do capítulo}
\addcontentsline{toc}{section}{Resumo do capítulo}
\begin{itemize}
\item $f'(a)=\lim_{h\to0}\frac{f(a+h)-f(a)}h$: inclinação da tangente e taxa de variação
  instantânea. Derivável implica contínua, mas não o contrário (bicos, tangentes verticais,
  oscilações).
\item Regras: $(x^r)'=rx^{r-1}$; $(f\pm g)'=f'\pm g'$; $(fg)'=f'g+fg'$;
  $(f/g)'=(f'g-fg')/g^2$; $(f\circ g)'=f'(g)\,g'$.
\item Trigonométricas: $(\sin x)'=\cos x$, $(\cos x)'=-\sin x$, $(\tan x)'=\sec^2x$,
  $(\sec x)'=\sec x\tan x$.
\item Inversas: $(f^{-1})'(x)=1/f'(f^{-1}(x))$; $(\sin^{-1}x)'=1/\sqrt{1-x^2}$,
  $(\tan^{-1}x)'=1/(1+x^2)$.
\item Exponenciais e logaritmos: $(e^x)'=e^x$, $(\ln x)'=1/x$, $(b^x)'=b^x\ln b$.
\item Movimento: $v=s'$, $a=v'=s''$. Economia: custo, receita e lucro marginais são
  derivadas. Variação aproximada: $f(a+h)\approx f(a)+f'(a)h$.
\item Derivação implícita: derive tratando $y$ como função de $x$ e isole $y'$.
  Derivação logarítmica para $f(x)^{g(x)}$ e produtos complicados.
\end{itemize}

% Gerado por ferramentas/gerar.py; não edite à mão.
\section*{Exercícios}\markright{Exercícios}
\addcontentsline{toc}{section}{Exercícios}
Os pontos antes de cada número indicam a dificuldade: \nivel{1} fácil, \nivel{2} médio, \nivel{3} difícil e \nivel{4} desafio. Os exercícios marcados com \estrela{} têm solução completa, com animação do gráfico, no repositório do livro: \url{https://github.com/AleffAzevedo/calculo-1-pt/blob/main/exercicios/cap3/README.md}. As respostas estão no fim do capítulo.

\begin{itemize}[leftmargin=5.2em,labelsep=0.6em,itemsep=4pt]
\item[\exnum{1}{3.1}] Use a definição de derivada para calcular $f'(x)$ para $f(x)=x^2-3x$.
\item[\exnum{1}{3.2}] Use a definição para calcular $f'(x)$ para $f(x)=\dfrac1{x+1}$.
\item[\exnum{1}{3.3}] Use a definição para calcular $f'(x)$ para $f(x)=\sqrt{2x+1}$. (Dica: conjugado.)
\item[\exnum{1}{3.4}] Encontre a reta tangente a $y=x^3$ em $x=2$.
\item[\exnum{1}{3.5}] Derive $y=4x^5-\dfrac3{x^2}+2\sqrt x-7$.
\item[\exnum{1}{3.6}] Derive $y=(x^2+1)(x^3-2x)$ de duas formas: pela regra do produto e expandindo antes.
\item[\exnum{1}{3.7}] Derive $y=\dfrac{2x-1}{x+3}$.
\item[\exnum{1}{3.8}] Derive $y=\dfrac{x^2}{x^2+1}$.
\item[\exnum{1}{3.9}] Derive $y=x^2\sen x$.
\item[\exnum{1}{3.10}] Derive $y=\dfrac{\tg x}{x}$.
\item[\exnum{1}{3.11}] Derive $y=(3x^2-5)^4$.
\item[\exnum{1}{3.12}] Derive $y=\sen(4x+1)$ e $y=\cos^2x$.
\item[\exnum{1}{3.13}] Derive $y=\sqrt{1+x^2}$.
\item[\exnum{1}{3.14}] Derive $y=e^{-3x}\cos x$.
\item[\exnum{1}{3.15}] Derive $y=\ln(x^2+4)$ e $y=\ln(\sen x)$.
\item[\exnum{1}{3.16}] Derive $y=x^3\,2^x$.
\item[\exnum{1}{3.17}] Derive $y=\arcsen(2x)$ e $y=\arctg(3x)$.
\item[\exnum{1}{3.18}] Calcule $y''$ para $y=x^4-2x^3+x$ e para $y=\sen(2x)$.
\item[\exnum{1}{3.19}] Para $s(t)=t^3-6t^2+9t$ (m), calcule a velocidade e a aceleração em $t=2$ s.
\item[\exnum{1}{3.20}] Use derivação implícita para encontrar $y'$ em $x^2+xy+y^2=7$.
\item[\exnum{2}{3.21}] \estrela\ \textbf{Retas tangente e normal.} Encontre as retas tangente e normal à parábola $y=x^2-4x+5$ no ponto de abscissa
$x=3$ e esboce as três curvas.
\item[\exnum{2}{3.22}] \estrela\ \textbf{Construindo o gráfico de $f'$.} Para $f(x)=x^3-3x^2$, determine onde $f'(x)>0$, $f'(x)<0$ e $f'(x)=0$, e esboce
o gráfico de $f'$ alinhado com o de $f$.
\item[\exnum{2}{3.23}] \estrela\ \textbf{Função por partes derivável.} Encontre $a$ e $b$ para que $f(x)=\begin{cases}x^2,&x\le1\\ax+b,&x>1\end{cases}$
seja derivável em todo $\mathbb R$.
\item[\exnum{2}{3.24}] Derive $y=(x^2+1)^3(2x-1)^2$ e fatore o resultado.
\item[\exnum{2}{3.25}] Derive $y=\sen^2(3x)$ e simplifique.
\item[\exnum{2}{3.26}] Derive $y=\cos\sqrt x$.
\item[\exnum{2}{3.27}] Derive $y=x\,e^{-x^2}$ e encontre onde $y'=0$.
\item[\exnum{2}{3.28}] Derive $y=\ln(\cos x)$ e $y=\ln|\sec x+\tg x|$.
\item[\exnum{2}{3.29}] Derive $y=\arctg\sqrt x$.
\item[\exnum{2}{3.30}] \estrela\ \textbf{O quociente de diferenças do seno.} Para $h=1$, $0{,}5$ e $0{,}1$, esboce $q_h(x)=\dfrac{\sen(x+h)-\sen x}{h}$ junto com
$\cos x$. O que acontece quando $h\to0$? Prove o resultado.
\item[\exnum{2}{3.31}] \estrela\ \textbf{Regra da cadeia e tangente.} Encontre a reta tangente a $y=\sen(x^2)$ em $x=\sqrt{\pi}$.
\item[\exnum{2}{3.32}] Calcule $y''$ para $y=xe^x$ e deduza uma fórmula para $y^{(n)}$.
\item[\exnum{2}{3.33}] Calcule $\dfrac{d^{50}}{dx^{50}}\cos x$ e $\dfrac{d^{99}}{dx^{99}}\sen x$.
\item[\exnum{2}{3.34}] Encontre os pontos em que $y=\dfrac{x}{x^2+1}$ tem tangente horizontal.
\item[\exnum{2}{3.35}] Use $f(a+h)\approx f(a)+f'(a)h$ para estimar $\sqrt{4{,}1}$ e $\sqrt[3]{8{,}3}$.
\item[\exnum{2}{3.36}] Sabendo que $g(1)=2$, $g'(1)=3$, $f(2)=5$ e $f'(2)=-1$, calcule $h'(1)$ para $h=f\circ g$ e para $k=g\cdot(f\circ g)$.
\item[\exnum{2}{3.37}] \estrela\ \textbf{Derivada de uma inversa sem a fórmula.} Seja $f(x)=x+e^x$. Mostre que $f$ é inversível e calcule $(f^{-1})'(1)$ e a reta
tangente ao gráfico de $f^{-1}$ no ponto $(1,0)$.
\item[\exnum{2}{3.38}] Use as propriedades do logaritmo antes de derivar: $y=\ln\dfrac{x^2+1}{x-1}$.
\item[\exnum{2}{3.39}] \estrela\ \textbf{Custo marginal.} O custo (em mil reais) de produzir $x$ unidades é $C(x)=0{,}01x^3-0{,}6x^2+15x+200$.
Calcule o custo marginal em $x=20$ e compare com o custo real de produzir a 21ª unidade.
\item[\exnum{2}{3.40}] Derive $y=3^{\sen x}$ e $y=\log_2(x^2+1)$.
\item[\exnum{3}{3.41}] \estrela\ \textbf{Partícula em movimento.} Uma partícula tem posição $s(t)=t^3-9t^2+24t$ (metros, $t\ge0$ em segundos).
(a) Quando ela está em repouso? (b) Quando se move para a frente e para trás?
(c) Qual é a distância total percorrida em $[0,5]$? (d) Quando está acelerando?
\item[\exnum{3}{3.42}] \estrela\ \textbf{Tangentes horizontais numa elipse.} Encontre os pontos da curva $x^2-xy+y^2=3$ onde a reta tangente é horizontal e onde é
vertical.
\item[\exnum{3}{3.43}] \estrela\ \textbf{O mínimo de $x^x$.} Use derivação logarítmica para encontrar $\dfrac{d}{dx}x^x$ ($x>0$) e determine o ponto onde
o gráfico de $y=x^x$ tem tangente horizontal.
\item[\exnum{3}{3.44}] \estrela\ \textbf{Tangentes por um ponto externo.} Encontre todas as retas tangentes à parábola $y=x^2$ que passam pelo ponto $(1,-3)$,
que não está na parábola.
\item[\exnum{3}{3.45}] Prove, pela definição, que $\dfrac{d}{dx}\cos x=-\sen x$.
\item[\exnum{3}{3.46}] Deduza $\dfrac{d}{dx}\arctg x=\dfrac1{1+x^2}$ a partir do teorema da derivada da inversa.
\item[\exnum{3}{3.47}] Derive $y=x^{\sen x}$ ($x>0$).
\item[\exnum{3}{3.48}] Use derivação logarítmica: $y=\dfrac{(x+1)^2(x-2)^3}{(x+3)^4}$.
\item[\exnum{3}{3.49}] Para a circunferência $x^2+y^2=r^2$, mostre que $y''=-\dfrac{r^2}{y^3}$.
\item[\exnum{3}{3.50}] Encontre a reta tangente à curva $x^2y+xy^2=6$ no ponto $(1,2)$.
\item[\exnum{3}{3.51}] Onde $f(x)=|x^2-4|$ é derivável? Calcule as derivadas laterais nos pontos problemáticos.
\item[\exnum{3}{3.52}] Mostre que $(\senh x)'=\cosh x$, $(\cosh x)'=\senh x$ e $(\tgh x)'=\operatorname{sech}^2x$.
\item[\exnum{3}{3.53}] Seja $f(x)=x^2\sen\frac1x$ para $x\neq0$ e $f(0)=0$. Calcule $f'(0)$ pela definição e mostre que $f'$ não é contínua em $0$.
\item[\exnum{3}{3.54}] A reta normal a $y=e^x$ em $x=0$ corta o eixo $x$ em que ponto?
\item[\exnum{3}{3.55}] Para que valor de $a$ a reta $y=ax$ é tangente à curva $y=e^x$? Em que ponto?
\item[\exnum{3}{3.56}] Derive $y=\arcsen x+\arccos x$ e conclua uma identidade.
\item[\exnum{3}{3.57}] Mostre que $y=e^{2x}+e^{-3x}$ satisfaz $y''+y'-6y=0$.
\item[\exnum{3}{3.58}] Encontre a reta tangente a $y=x^x$ em $x=e$.
\item[\exnum{3}{3.59}] O ângulo de um pêndulo é $\theta(t)=0{,}2\cos(2\pi t)$ rad. Qual é a velocidade angular máxima e quando ocorre?
\item[\exnum{3}{3.60}] Em que pontos a tangente a $y=x^3$ é paralela à reta $y=12x+1$? Escreva as tangentes.
\item[\exnum{4}{3.61}] Prove que $(fgh)'=f'gh+fg'h+fgh'$ e generalize para $n$ fatores.
\item[\exnum{4}{3.62}] Para que valores de $n$ a função $y=x^n$ satisfaz $x^2y''-3xy'+4y=0$?
\item[\exnum{4}{3.63}] Prove por indução que $\dfrac{d^n}{dx^n}(xe^x)=(x+n)e^x$.
\item[\exnum{4}{3.64}] Mostre que as curvas $x^2-y^2=a$ e $xy=b$ ($a,b\neq0$) se cortam em ângulo reto.
\item[\exnum{4}{3.65}] Encontre todas as retas que são tangentes simultaneamente a $y=x^2$ e a $y=-x^2+4x-4$.
\end{itemize}
\subsection*{Respostas}\markright{Respostas}
{\small\begin{multicols}{2}\raggedright
\begin{itemize}[leftmargin=2.7em,labelsep=0.5em,itemsep=1pt]
\item[\textbf{3.1}] $2x-3$.
\item[\textbf{3.2}] $-\dfrac1{(x+1)^2}$.
\item[\textbf{3.3}] $\dfrac1{\sqrt{2x+1}}$.
\item[\textbf{3.4}] $y=12x-16$.
\item[\textbf{3.5}] $20x^4+\dfrac6{x^3}+\dfrac1{\sqrt x}$.
\item[\textbf{3.6}] $5x^4-3x^2-2$.
\item[\textbf{3.7}] $\dfrac7{(x+3)^2}$.
\item[\textbf{3.8}] $\dfrac{2x}{(x^2+1)^2}$.
\item[\textbf{3.9}] $2x\sen x+x^2\cos x$.
\item[\textbf{3.10}] $\dfrac{x\sec^2x-\tg x}{x^2}$.
\item[\textbf{3.11}] $24x(3x^2-5)^3$.
\item[\textbf{3.12}] $4\cos(4x+1)$; $-2\sen x\cos x=-\sen2x$.
\item[\textbf{3.13}] $\dfrac{x}{\sqrt{1+x^2}}$.
\item[\textbf{3.14}] $-e^{-3x}(3\cos x+\sen x)$.
\item[\textbf{3.15}] $\dfrac{2x}{x^2+4}$; $\cotg x$.
\item[\textbf{3.16}] $2^x\big(3x^2+x^3\ln2\big)$.
\item[\textbf{3.17}] $\dfrac2{\sqrt{1-4x^2}}$; $\dfrac3{1+9x^2}$.
\item[\textbf{3.18}] $12x^2-12x$; $-4\sen2x$.
\item[\textbf{3.19}] $v(2)=-3$ m/s; $a(2)=0$.
\item[\textbf{3.20}] $y'=-\dfrac{2x+y}{x+2y}$.
\item[\textbf{3.21}] Tangente $y=2x-4$; normal $y=-\tfrac12x+\tfrac72$.
\item[\textbf{3.22}] $f'(x)=3x(x-2)$: positiva em $(-\infty,0)\cup(2,\infty)$, negativa em $(0,2)$, nula em $x=0$ e $x=2$.
\item[\textbf{3.23}] $a=2$, $b=-1$.
\item[\textbf{3.24}] $2(x^2+1)^2(2x-1)(8x^2-3x+2)$.
\item[\textbf{3.25}] $6\sen3x\cos3x=3\sen6x$.
\item[\textbf{3.26}] $-\dfrac{\sen\sqrt x}{2\sqrt x}$.
\item[\textbf{3.27}] $e^{-x^2}(1-2x^2)$; $x=\pm\tfrac1{\sqrt2}$.
\item[\textbf{3.28}] $-\tg x$; $\sec x$.
\item[\textbf{3.29}] $\dfrac1{2\sqrt x\,(1+x)}$.
\item[\textbf{3.30}] $q_h\to\cos x$; usa $\lim\frac{\sen h}h=1$ e $\lim\frac{\cos h-1}h=0$.
\item[\textbf{3.31}] $y=-2\sqrt\pi\,(x-\sqrt\pi)$.
\item[\textbf{3.32}] $y''=(x+2)e^x$; $y^{(n)}=(x+n)e^x$.
\item[\textbf{3.33}] $-\cos x$; $-\cos x$.
\item[\textbf{3.34}] $(1,\tfrac12)$ e $(-1,-\tfrac12)$.
\item[\textbf{3.35}] $2{,}025$; $2{,}025$.
\item[\textbf{3.36}] $h'(1)=-3$; $k'(1)=3\cdot5+2\cdot(-3)=9$.
\item[\textbf{3.37}] $(f^{-1})'(1)=\tfrac12$; tangente $y=\tfrac12(x-1)$.
\item[\textbf{3.38}] $\dfrac{2x}{x^2+1}-\dfrac1{x-1}$.
\item[\textbf{3.39}] $C'(20)=3$; o custo real é $C(21)-C(20)=3{,}01$.
\item[\textbf{3.40}] $3^{\sen x}\ln3\cos x$; $\dfrac{2x}{(x^2+1)\ln2}$.
\item[\textbf{3.41}] (a) $t=2$ e $t=4$; (b) frente em $[0,2)\cup(4,\infty)$, trás em $(2,4)$; (c) 28 m; (d) em $(2,3)$ e $(4,\infty)$.
\item[\textbf{3.42}] Horizontal em $(1,2)$ e $(-1,-2)$; vertical em $(2,1)$ e $(-2,-1)$.
\item[\textbf{3.43}] $y'=x^x(\ln x+1)$; tangente horizontal em $x=1/e$, com $y=e^{-1/e}\approx0{,}692$.
\item[\textbf{3.44}] $y=6x-9$ (tangente em $x=3$) e $y=-2x-1$ (tangente em $x=-1$).
\item[\textbf{3.45}] Use $\cos(x+h)=\cos x\cos h-\sen x\sen h$ e os limites do Teorema 2.7.
\item[\textbf{3.46}] $y=\arctg x\Rightarrow y'=\dfrac1{\sec^2y}=\dfrac1{1+\tg^2y}$.
\item[\textbf{3.47}] $x^{\sen x}\Big(\cos x\ln x+\dfrac{\sen x}x\Big)$.
\item[\textbf{3.48}] $y\Big(\dfrac2{x+1}+\dfrac3{x-2}-\dfrac4{x+3}\Big)$.
\item[\textbf{3.49}] $y'=-\tfrac xy$; derive de novo e use $x^2+y^2=r^2$.
\item[\textbf{3.50}] $y'=-\tfrac85$; $y=-\tfrac85x+\tfrac{18}5$.
\item[\textbf{3.51}] Em todo $x\neq\pm2$; em $x=2$: $-4$ e $4$; em $x=-2$: $-4$ e $4$.
\item[\textbf{3.52}] Derive as definições com exponenciais; use $\cosh^2-\senh^2=1$.
\item[\textbf{3.53}] $f'(0)=0$; $f'(x)=2x\sen\tfrac1x-\cos\tfrac1x$ não tem limite em 0.
\item[\textbf{3.54}] Normal $y=-x+1$; corta em $(1,0)$.
\item[\textbf{3.55}] $a=e$, no ponto $(1,e)$.
\item[\textbf{3.56}] $y'=0$, logo $\arcsen x+\arccos x=\tfrac\pi2$.
\item[\textbf{3.57}] Substitua $y'$ e $y''$; os termos se cancelam.
\item[\textbf{3.58}] $y=e^e+2e^e(x-e)$.
\item[\textbf{3.59}] $0{,}4\pi\approx1{,}26$ rad/s, quando $\theta=0$ ($t=\tfrac14+\tfrac n2$).
\item[\textbf{3.60}] $x=\pm2$: $y=12x-16$ e $y=12x+16$.
\item[\textbf{3.61}] Aplique a regra do produto a $(fg)\cdot h$; por indução, soma de $n$ termos com uma derivada em cada.
\item[\textbf{3.62}] $n=2$ (raiz dupla de $n^2-4n+4=0$).
\item[\textbf{3.63}] Base $n=1$; passo: derive $(x+n)e^x$ e obtenha $(x+n+1)e^x$.
\item[\textbf{3.64}] Inclinações $\tfrac xy$ e $-\tfrac yx$: o produto é $-1$.
\item[\textbf{3.65}] $y=0$ e $y=4x-4$.
\end{itemize}\end{multicols}}

